%! AP Statistics — Unit 2, Lesson 2.3 — Student Packet Cover Sheet
\documentclass[10pt]{article}
\usepackage{apstats-article}
\usepackage{apstats-boxes}

%=============================================================================
\begin{document}

% ── Full-bleed navy banner ────────────────────────────────────────────────────
\begin{tikzpicture}[remember picture, overlay]
  \fill[navy] (current page.north west)
    rectangle ([yshift=-0.9in]current page.north east);
\end{tikzpicture}
\vspace{-0.8in}

\begin{center}
  {\color{white}\textbf{\LARGE AP Statistics}}\\[4pt]
  {\color{skymid}\large\textbf{Unit 2: Exploring Two-Variable Data}}\\[6pt]
  {\color{white}\textbf{\large Lesson 2.3 \quad Statistics for Two Categorical Variables}}
\end{center}

\vspace{0.22in}

\namedateperiod

\vspace{0.18in}

% ── LEARNING TARGETS ─────────────────────────────────────────────────────────
\begin{learningtargetbox}
\small
\begin{enumerate}[leftmargin=1.5em, itemsep=5pt]
  \item \ldots{} compute the conditional proportion (relative frequency) for each group in a
        two-way table. \hfill\textcolor{slate}{\small(UNC-1.Q)}
  \item \ldots{} calculate and interpret the \textbf{difference in proportions} between two
        groups in context. \hfill\textcolor{slate}{\small(UNC-1.Q, UNC-1.R)}
  \item \ldots{} calculate and interpret \textbf{relative risk} (the ratio of two conditional
        proportions) in context. \hfill\textcolor{slate}{\small(UNC-1.Q, UNC-1.R)}
  \item \ldots{} compare the difference in proportions and relative risk to determine which
        better describes the strength of an association. \hfill\textcolor{slate}{\small(UNC-1.R)}
\end{enumerate}
\end{learningtargetbox}

\vspace{0.14in}

% ── TABLE OF CONTENTS ────────────────────────────────────────────────────────
\begin{tocbox}
\small
\renewcommand{\arraystretch}{1.5}
\begin{tabularx}{\linewidth}{@{} c l X r @{}}
\rowcolor{sky}
\textbf{\#} & \textbf{Component} & \textbf{Description} & \textbf{Score} \\
\midrule
1 & Warm-Up       & Spiral review --- conditional relative frequencies         & \blank{1.2cm} \\
2 & Guided Notes  & Difference in proportions and relative risk                 & \blank{1.2cm} \\
3 & Group Activity & Compute and compare both statistics from real-world data   & \blank{1.2cm} \\
4 & Exit Ticket   & Calculate and interpret difference in proportions \& RR     & \blank{1.2cm} \\
5 & Homework      & Multi-context practice with both summary statistics          & \blank{1.2cm} \\
\midrule
  & \multicolumn{2}{r}{\textbf{Total}} & \blank{1.2cm} \\
\end{tabularx}
\end{tocbox}

\vspace{0.14in}

% ── KEEP IN MIND ─────────────────────────────────────────────────────────────
\begin{remindbox}
\small
Two statistics quantify associations between categorical variables. Here are the key ideas.

\vspace{6pt}
\begin{tabularx}{\linewidth}{@{} >{\centering\arraybackslash\columncolor{goldbg}}p{0.06\linewidth}
                                    X @{}}
\textbf{\large!} &
\textbf{Always condition on the explanatory variable.}\enspace
Compute $\hat{p}$ for each group by dividing the success count by that group's total,
not the grand total. \\[6pt]
\textbf{\large!} &
\textbf{Difference in proportions is additive.}\enspace
$\hat{p}_1 - \hat{p}_2$ answers ``how many percentage points higher?''
A value of 0 means no difference; sign tells you direction. \\[6pt]
\textbf{\large!} &
\textbf{Relative risk is multiplicative.}\enspace
$\hat{p}_1 / \hat{p}_2$ answers ``how many times as likely?''
A value of 1 means no difference; values $> 1$ or $< 1$ signal direction. \\[6pt]
\textbf{\large!} &
\textbf{One number isn't enough.}\enspace
A small absolute difference can correspond to a large relative risk (and vice versa).
Always report and interpret both in context.
\end{tabularx}
\end{remindbox}

\end{document}
